\notechapter{N-20220716}{Conic Exercises: Tangency, Intersections at Infinity, and Polarity}

Substituting a line into a conic produces a quadratic equation, but a repeated root and a single \emph{finite} intersection are different statements. The distinction becomes clear in the projective description of conics, as developed in the Oxford projective geometry lectures.\footnote{\url{https://courses.maths.ox.ac.uk/course/section.php?id=10755}}

\section{A secant of an ellipse}

Consider the line \(y=2x-2\) and the ellipse
\[
 \frac{x^2}{5}+\frac{y^2}{4}=1.
\]
After substitution,
\[
 \frac{x^2}{5}+\frac{(2x-2)^2}{4}=1
 \quad\Longleftrightarrow\quad
 \frac65x^2-2x=0.
\]
The roots \(x=0\) and \(x=5/3\) give the intersection points
\[
 (0,-2),\qquad (5/3,4/3).
\]
Because the equation is genuinely quadratic, the discriminant distinguishes two real intersections, a repeated contact point, and no real intersection.

\section{Tangents to an ellipse from a specified point}

Let \(E:x^2/a^2+y^2/b^2=1\) with \(a,b>0\). For the point \(P=(0,2)\), write a possible tangent as \(y=kx+2\). Substitution gives
\[
 \left(\frac1{a^2}+\frac{k^2}{b^2}\right)x^2
 +\frac{4k}{b^2}x+\frac4{b^2}-1=0.
\]
The discriminant vanishes exactly when
\[
 \boxed{k^2=\frac{4-b^2}{a^2}}.
\]
Consequently there are two real tangents when \(b<2\), one horizontal tangent when \(b=2\), and no real tangents when \(b>2\). The particular numerical slopes depend on \emph{both} semiaxes. A claimed pair of slopes without specifying \(a,b\) is insufficiently determined.

The same result follows by using the tangent at \((x_0,y_0)\in E\),
\[
 \frac{x x_0}{a^2}+\frac{y y_0}{b^2}=1.
\]
It passes through \((0,2)\) precisely when \(y_0=b^2/2\). Substituting into the ellipse gives the same existence condition. This is a useful independent check of the discriminant calculation.

\section{One Intersection without Tangency}

For \(x^2-y^2=2\) and \(y=x-1\), substitution gives
\[
 x^2-(x-1)^2=2,\qquad x=\frac32,\quad y=\frac12.
\]
The resulting equation is linear, so its unique affine solution is not a double root. The projective hyperbola is \(X^2-Y^2=2Z^2\), and the line is \(Y=X-Z\). At infinity \(Z=0\), these equations share the point \([1:1:0]\). The missing second projective intersection has moved to infinity, along an asymptotic direction. This explains why merely counting finite points is an unsafe test for tangency.

\section{Correcting a parabola tangent calculation}

Take \(y=2x^2\) and \(P=(2,6)\). A line through \(P\) of slope \(k\) has equation \(y=kx+6-2k\). Contact requires a repeated root of
\[
 2x^2-kx-6+2k=0.
\]
The discriminant equation is
\[
 k^2-8(-6+2k)=k^2-16k+48=(k-4)(k-12)=0.
\]
Thus the actual tangents are
\[
 \boxed{y=4x-2,\qquad y=12x-18.}
\]
The previously written negative slope does not satisfy the discriminant and is not a tangent through the stated point. We can verify the new answers by substitution: the first yields \(2(x-1)^2=0\); the second yields \(2(x-3)^2=0\).

\section{A Focus-and-Circle Maximum}

Consider \(x^2=2py\) with \(p>0\), focus \(F=(0,p/2)\), and the circle
\[
 x^2+(y+4)^2=1.
\]
Its center is \((0,-4)\); because \(F\) lies above the circle, the shortest focus-to-circle distance is \(p/2+3\). If this distance is \(4\), then \(p=2\), giving the parabola \(x^2=4y\).

Parametrize tangent contact points by \(T(t)=(t,t^2/4)\). Its tangent equation is
\[
 y=\frac t2x-\frac{t^2}{4}.
\]
For an intersection point \(P=(u,v)\), the two contact parameters \(t_1,t_2\) solve
\[
 t^2-2ut+4v=0,\qquad
 t_1+t_2=2u,\quad t_1t_2=4v.
\]
They are real and distinct when \(u^2-4v>0\). Directly computing the determinant of the two vectors from \(P\) to \(T(t_1)\) and \(T(t_2)\) gives
\[
 [PT(t_1)T(t_2)]
 =\frac{|t_1-t_2|^3}{16}
 =\frac12(u^2-4v)^{3/2}.
\]
On the circle, \(u^2+(v+4)^2=1\), where \(-5\le v\le-3\). Therefore
\[
 u^2-4v=1-(v+4)^2-4v=21-(v+6)^2\le20.
\]
Equality is attained at \(P=(0,-5)\). Hence
\[
 \boxed{[PAB]_{\max}=\tfrac12 20^{3/2}=20\sqrt5.}
\]
The value comes from a verifiable extremum and explicit tangent geometry, not from an unexplained final number.


\section{An ellipse area constraint requires specified vertices}

Suppose the quadrilateral is formed by the four axis-intercepts
\[
 (\pm a,0),\qquad (0,\pm b)
\]
of \(x^2/a^2+y^2/b^2=1\). The quadrilateral is a rhombus with diagonals \(2a\) and \(2b\), so its area is \(2ab\). If it passes through \((0,-2)\) and the rhombus area is \(4\sqrt5\), then \(b=2\), \(a=\sqrt5\), and the ellipse is
\[
 \frac{x^2}{5}+\frac{y^2}{4}=1.
\]

\begin{caution}
The condition that an unspecified quadrilateral inscribed in an ellipse has area \(4\sqrt5\) does not determine \(a\) from \(b=2\). The deduction \(2ab=4\sqrt5\) requires the stated axial vertices or an equivalent geometric condition.
\end{caution}

For the line \(y=kx-3\) through \((0,-3)\), substituting into this ellipse produces
\[
 (4+5k^2)x^2-30kx+25=0.
\]
Its discriminant is
\[
 (-30k)^2-100(4+5k^2)=400(k^2-1).
\]
Thus two distinct real intersections occur exactly when \(|k|>1\), tangency when \(|k|=1\), and no real intersections when \(|k|<1\). Additional length constraints require the missing geometric definitions of their endpoints; no range for \(k\) should be quoted without them.

\section{Polarity packages tangent equations}

Over a field of characteristic different from two, let a nonsingular projective conic be \(X^TAX=0\), with \(A\) symmetric and invertible. At a point \(T\) of the conic, its tangent line is
\[
 T^TAX=0.
\]
This follows by differentiating the quadratic equation:
\[
 d(X^TAX)_T(H)=2T^TAH.
\]
For an arbitrary projective point \(P\), the line \(P^TAX=0\) is its \emph{polar}. If \(P\) lies outside a real conic and has two real tangent-contact points \(T_1,T_2\), then those points both satisfy \(P^TAT_i=0\), so the polar is the chord joining them.

The projective argument clarifies why tangent-chord relations can be handled independently of Euclidean distances. The discriminant counts multiplicity on a line; the bilinear form describes tangency; and affine normalization can simplify ellipse calculations when an area factor is tracked by the determinant of the change of variables. None of these tools requires replacing an actual calculation with a slogan.
